\subsection{Earlier bounds for Fourier circuits}

The classical $O(N\log N)$ count has two ingredients. Cooley and
Tukey factor a transform of length $ab$ into transforms along the two
axes, with a diagonal scaling between them. Good's factorization
eliminates the scaling for coprime lengths
\cite{Good,CooleyTukey,cooleyLewisWelch1967}. Recursing on powers of
two gives the familiar FFT. Bluestein's reduction to convolution
with a fixed quadratic phase, followed by padding to a power of two,
gives the same asymptotic count for every complex transform length
\cite{Bluestein,rabinerSchaferRader1969}.

Several classical lower bounds explain when this count is necessary.
Papadimitriou proved optimality under restrictions on the paths in
the computation graph \cite{Papadimitriou}. Winograd counted multiplications
separately from additions \cite{Winograd1979DFT}. For the unnormalized
complex transform, the determinant has modulus $N^{N/2}$.
Morgenstern showed that if scalar multipliers have modulus at most one,
an addition can at most double the largest determinant modulus
attainable from the available forms. At least $\tfrac12N\log_2 N$
additions are therefore necessary \cite{Morgenstern}. Normalizing the
transform makes its determinant have modulus one. For this transform,
Ailon proved an $\Omega(N\log N)$ lower bound for in-place
two-coordinate unitary operations, and an $\Omega(N\log N/R)$ bound
for in-place rotations and scalings when every intermediate map has
condition number at most $R$ \cite{Ailon,AilonCondition}.
Theorem~\ref{thm:main} permits arbitrary scalar coefficients,
intermediate storage and reuse.

For power-of-two lengths, Alman and Rao used matrix nonrigidity to
reduce the leading operation counts to $\tfrac{23}{24}N\log_2 N$
for Walsh--Hadamard transforms over a field and $\tfrac{15}{4}N\log_2 N$
in real arithmetic for complex cyclic transforms, up to $O(N)$ terms
\cite{AlmanRao}. Finite tensor constructions subsequently gave
$o(N\log N)$ circuits along a subsequence \cite{FiniteFourier}.
At every length, the algorithm of \cite[Theorem~1.1]{ExplicitFourier}
uses a finite network from \cite[Proposition~8 and Section~3.5]{Tape}
to save a fixed positive power of $\log N$. It charges for array indexing and for preparing the
scalars from a supplied root of unity of polynomial order.
Sussman developed a recursion that handles several target maps
together \cite{Sussman}. Kartal combined it with constructions of
CubeJerry, Cinna Davis and Ali Romman, and Yao and Liu refined the
resulting numerical bounds \cite{Kartal,kartalpi,YaoLiu}. A subsequent
construction gives a saving greater than $0.0039$ in this complex
arithmetic model \cite{RelayFourier}. The circuit size in Theorem~\ref{thm:main}
charges only for the binary linear gates.

Valiant related linear circuits to matrix rigidity, the number of
entries that must change to reach a prescribed rank
\cite{Valiant1975,Valiant1977}. Alman and Williams showed that the
Walsh matrix fails Valiant's rigidity criterion \cite{AlmanWilliamsRigidity}.
For every fixed $\delta>0$, Dvir and Liu reduce the rank of any
sufficiently large complex abelian character table to
$o(N/\log\log N)$ by changing at most $N^{1+\delta}$ entries
\cite[Theorem~6.1]{DvirLiu}. Theorem~\ref{bch:uniform-robust},
with one transform and no circuit states, gives a lower bound in the
same variables. Reducing the rank to $\beta N$ requires more than
$\eta N2^{\sqrt{\log N}/8}$ changes. For a smaller target rank, Alman
showed that reducing a Walsh matrix over $\mathbb F_3$ to rank
$\lfloor(\log_2 N)^2/80\rfloor$ requires changing more than one percent
of its entries \cite{AlmanRigidity}.

For convolution, specializing one input turns a bilinear circuit into
a linear one. Raz used specialization for bounded-coefficient matrix
multiplication \cite{Raz}. B\"urgisser and Lotz proved an
$\Omega(n\log n)$ bound for complex cyclic convolution and ordinary
polynomial multiplication in the bounded-coefficient bilinear model
\cite{BurgisserLotz}.
Corollary~\ref{cor:conv-exponent} uses a fixed operand over the algebraic
closure to obtain a Fourier transform, giving a lower bound for
division-free formal polynomial multiplication with arbitrary scalars
over every field.
